% =============================================================================
% SOMMAIRE DU CHAPITRE 1 — version mise à jour, 2 avril 2026
% Dix sections. Titres purement techniques.
% =============================================================================

\chapter{Transmettre, stocker, calculer~: l'émergence des techniques de l'information au regard de l'économie et de l'énergie}\label{ch:history}

\subsection*{Méthode, protocole et grille de lecture}\label{subsec:grammaire}


\section{Avant l'électricité (avant 1794)}\label{sec:era0}

\subsection{L'écriture comptable et les supports de mémoire}\label{subsec:ecriture}

\subsection{Le transport de l'information}\label{subsec:transport}

\subsection{L'imprimerie et la reproduction du savoir}\label{subsec:imprimerie}

\subsection{Pascal, Leibniz et les premières machines à calculer}\label{subsec:pascaline}

\subsection{Prony, le Nautical Almanac et la division du travail intellectuel}\label{subsec:prony}

\subsection{Le métier Jacquard et le contrôle par l'information stockée}\label{subsec:jacquard}


\section{$\blacktriangleright$~Le télégraphe (1794--1880)}\label{sec:telegraph}

\subsection{Le télégraphe optique et électrique}\label{subsec:telegraph}

\subsection{$\blacktriangleright$~Le télégraphe dans le chemin de fer}\label{subsec:railroad}


\section{La mécanisation du calcul (1820--1890)}\label{sec:calculating}

\subsection{L'Arithmomètre de Thomas de Colmar}\label{subsec:arithmometre}

\subsection{Babbage~: le Difference Engine et l'Analytical Engine}\label{subsec:babbage}

\subsection{Burroughs, Odhner, Brunsviga et la caisse enregistreuse NCR}\label{subsec:burroughs}

\subsection{$\blacktriangleright$~Hollerith, le recensement de 1890 et la naissance de la tabulatrice}\label{subsec:hollerith}


\section{Le téléphone, la radio et les débuts de l'enregistrement (1876--1920)}\label{sec:telephone}

\subsection{Bell, Gray et l'invention du téléphone}\label{subsec:bell}

\subsection{AT\&T et la commutation}\label{subsec:att}

\subsection{Marconi et la télégraphie sans fil}\label{subsec:marconi}

\subsection{Fleming, De Forest et le tube à vide}\label{subsec:tube}

\subsection{Le phonographe, le gramophone et le cinématographe}\label{subsec:phonographe}


\section{L'électrification, les tabulatrices et la radiodiffusion (1920--1945)}\label{sec:electrification}

\subsection{$\blacktriangleright$~Le réseau électrique~: Edison, Insull et les centrales}\label{subsec:edison}

\subsection{La radio et la télévision de masse}\label{subsec:broadcasting}

\subsection{$\blacktriangleright$~IBM et les tabulatrices dans les grandes organisations}\label{subsec:ibm}

\subsection{Boole, Shannon et la logique devenue circuit}\label{subsec:shannon}

\subsection{Enigma, Bletchley Park et Colossus}\label{subsec:enigma}


\section{Le calcul programmable et la révolution électronique (1945--1975)}\label{sec:computing}

\subsection{L'ENIAC et l'architecture de von Neumann}\label{subsec:eniac}

\subsection{IBM et les grandes machines}\label{subsec:mainframe}

\subsection{Le transistor, le circuit intégré et le microprocesseur}\label{subsec:transistor}

\subsection{ARPANET et la commutation par paquets}\label{subsec:arpanet}

\subsection{La concomitance de 1970--1975}\label{subsec:bascule}


\section{La convergence sur le réseau (1975--2007)}\label{sec:convergence}

\subsection{L'ordinateur personnel}\label{subsec:pc}

\subsection{Internet et le World Wide Web}\label{subsec:www}

\subsection{La bulle spéculative}\label{subsec:dotcom}

\subsection{Le téléphone intelligent}\label{subsec:smartphone}


\section{Paradigmes contemporains~: la convergence de la transmission, du stockage et du calcul (2007--)}\label{sec:contemporain}

\subsection{L'informatique en nuage et le centre de données à très grande échelle}\label{subsec:cloud}

\subsection{$\blacktriangleright$~Le stockage numérique et la mémoire permanente}\label{subsec:store}

\subsection{$\blacktriangleright$~L'intelligence artificielle générative et les lois d'échelle}\label{subsec:genai}

\subsection{La blockchain et les cryptomonnaies}\label{subsec:blockchain}

\subsection{L'IA agentique et l'automatisation cognitive}\label{subsec:agentique}

\subsection{L'Internet des objets et l'\emph{edge computing}}\label{subsec:iot}

\subsection{L'IA physique, la robotique et les véhicules autonomes}\label{subsec:physical_ai}

\subsection{Les jumeaux numériques et la simulation à grande échelle}\label{subsec:digital_twins}

\subsection{Les réseaux 5G/6G et l'infrastructure de transmission}\label{subsec:5g}


\section{Horizons~: paradigmes émergents et limites physiques}\label{sec:horizons}

\subsection{La limite thermodynamique~: calcul réversible et principe de Landauer}\label{subsec:reversible}

\subsection{Le retour au stockage passif~: ADN et mémoire sans énergie}\label{subsec:adn}

\subsection{Vers l'efficacité biologique~: calcul neuromorphique et photonique}\label{subsec:neuromorphique}

\subsection{Déplacer le calcul~: centres de données orbitaux, calcul quantique et intelligence artificielle générale}\label{subsec:deplacer}


\section{Des cas aux archétypes~: de l'abduction aux dimensions}\label{sec:reduction}